% uid: ""
% school: "lycee"
% level: "3eme"
% section: ["Economie et Gestion"]
% chapter: "Suites réelles"
% title: "Gestion d'un stock : suite récurrente"
% difficulty: "Moyen"
% points: null
% tags: ["suites", "suite-recurrente", "suite-auxiliaire", "suite-geometrique", "limite"]
% provenance: "own"
% dressing: "U(n+1)=aU(n)+b, a=0,75"
% rung: 3
% batch: "2026-11-20-suites-3eco-A"

% ---STATEMENT---
Un entrepôt contient $4\,000$ articles au début du mois $0$. Chaque mois, $25\,\%$ du stock est vendu puis $2\,000$ articles sont livrés. On note $U_n$ le stock au début du mois $n$.

\begin{enumerate}
  \item Montrer que $U_{n+1}=0{,}75\,U_n+2\,000$. Calculer $U_1$ et $U_2$.
  \item On pose $V_n=U_n-8\,000$. Montrer que $(V_n)$ est géométrique ; préciser sa raison et son premier terme.
  \item En déduire $U_n$ en fonction de $n$.
  \item Déterminer la limite de $(U_n)$. Interpréter.
\end{enumerate}

% ---SOLUTION---
\begin{enumerate}
  \item Après la vente de $25\,\%$ du stock il reste $75\,\%$, soit $0{,}75\,U_n$ ; puis on ajoute $2\,000$ articles :
  $$U_{n+1}=0{,}75\,U_n+2\,000.$$
  $$U_1=0{,}75\times4\,000+2\,000=5\,000\qquad U_2=0{,}75\times5\,000+2\,000=5\,750.$$
  \item \textbf{D'où vient $8\,000$ ?} C'est la valeur $\alpha$ pour laquelle le stock ne bougerait pas : $\alpha=0{,}75\,\alpha+2\,000\iff0{,}25\,\alpha=2\,000\iff\alpha=8\,000$.

  Calculons $V_{n+1}$ avec $V_n=U_n-8\,000$ :
  $$V_{n+1}=U_{n+1}-8\,000=0{,}75\,U_n+2\,000-8\,000=0{,}75\,U_n-6\,000=0{,}75\,(U_n-8\,000)=0{,}75\,V_n.$$
  $(V_n)$ est géométrique de raison $\ib{$0{,}75$}$ et de premier terme $V_0=4\,000-8\,000=$ \ib{$-4\,000$}.
  \item $V_n=-4\,000\times0{,}75^n$, donc $U_n=V_n+8\,000$ :
  $$ \ib{$U_n=8\,000-4\,000\times0{,}75^n$}.$$
  \emph{Vérification} : $U_1=8\,000-3\,000=5\,000$ $\checkmark$ \quad $U_2=8\,000-4\,000\times0{,}5625=5\,750$ $\checkmark$
  \item $0<0{,}75<1$, donc $0{,}75^n\to0$ et $U_n\to8\,000$.

  À long terme, le stock se stabilise à $8\,000$ articles (la valeur d'équilibre $\alpha$) : chaque mois, ce qu'on vend ($25\,\%$ de $8\,000=2\,000$) est exactement ce qu'on livre.
\end{enumerate}

\textbf{Idée :} le stock se rapproche de la valeur d'équilibre $\alpha$ où ce qui sort est égal à ce qui entre.
